

\section{Qu'est-ce que la recherche d'information (Information Retrieval)\,?}

\begin{frame}{\textit{Information Retrieval}, définition}

\begin{block}{Une définition}
La \alert{Recherche d'Information} (\textit{Information Retrieval}, IR) consiste à
trouver des \alert{document} peu ou faiblement structurés, dans une grande \alert{collection},
en fonction d'un \alert{besoin d'information}.
\end{block}

Prend peu à peu la prédominance sur les techniques de bases de données.

\begin{itemize} 
  \item Recherche sur le Web. Utilisée quotidiennement par des milliards d'utilisateurs.
   \item Recherche dans votre boîte mail.
   \item Recherche sur votre ordinateur (\textit{Spotlight}).
    \item Recherche dans une base documentaire, publique ou privée.
\end{itemize}

À bien distinguer d'un recherche type ``base de données''\,: requête structurée, données structurées.
\end{frame}

\begin{frame}{\textit{Information Retrieval}, scénario général}


\figSlideWithSize{ri-process}{11cm}

\end{frame}



\begin{frame}{Spécificités}

Recherche plus ou moins précise, documents peu structurés, résultats très nombreux. 

\begin{problem}
Comment évaluer la \alert{pertinence} du résultat\,?
\end{problem}

Deux critères esentiels (définition à préciser)\,:

\begin{itemize}
  \item \alert{Précision}\,: fraction des documents \alert{trouvés} qui sont pertinents pour l'utilisateur.
  \item \alert{Rappel}\,: fraction des documents \alert{pertinents} qui sont trouvés par le système.
\end{itemize}

\begin{block}{Réfléchissez-y\,!}
\begin{enumerate}
 \item quel est le rappel si la recherche ramène \alert{tous} les documents\,? 
 \item quelle est la précision si la recherche ne ramène \alert{aucun} document\,?
 \item et inversement...
 \end{enumerate}
 \end{block}
\end{frame}

\begin{frame}{Notre exemple de base}


\vspace{.5cm}

Un ensemble (modeste) de \alert{documents textuels}

\vspace{.5cm}

\begin{tabular}{>{\color{structure.fg}}ll}
$d_1$&The jaguar is a New World mammal of the Felidae family.\\
$d_2$&Jaguar has designed four new engines.\\
$d_3$&For Jaguar, Atari was keen to use a 68K family device.\\
$d_4$&The Jacksonville Jaguars are a professional US football team.\\
$d_5$&Mac OS X Jaguar is available at a price of US \$199 for
Apple's\\& new ``family pack''.\\
$d_6$&One such ruling family to incorporate the jaguar into their name\\&is
Jaguar Paw.\\
$d_7$&It is a big cat.
\end{tabular}

\medskip

Noter\,: les documents \alert{multimedia} soulèvent d'autres types de problèmes. Cf. NFE205.
\end{frame}


\subsection{Présentation informelle}


\begin{frame}{La Recherche d'Information au XVII$^e$ siècle}

\begin{block}{Mon besoin}
Dans quelle pièce de Shakespeare trouve-t-on Brutus, César, mais pas Calpurnia.
\end{block}
\vfill

Première solution\,: parcourir les textes des pièces exhaustivement.

\vfill

Pas très satisfaisant car\,:
\begin{itemize}
  \item c'est long,
  \item le critère ``pas Calpurnia'' n'est pas facile à traiter,
  \item autres types de recherche ("Brutus doit être \alert{près de} César") difficiles.
  \item \textit{quid} si je veux \alert{classer} par pertinence\,?
\end{itemize}

Les moteurs de recherche font tout cela\,!

\end{frame}

\begin{frame}{Un outil pour faire mieux et plus vite}

Une \alert{matrice d'incidence}.

\figSlideWithSize{incidence-matrix}{12cm}

\medskip

Parfois utilisée dans les bases de données sous le nom \alert{d'index bitmap}.
\end{frame}


\begin{frame}{Répondons à la recherche}

On prend les \alert{vecteurs d'incidence} de chaque terme.

\begin{itemize}
  \item Brutus:       110100
  \item César:        110111
  \item Calpurnia:    010000
\end{itemize}

On fait un \texttt{AND} sur les vecteurs de Brutus et César.

\vfill
\centerline{Résultat: 110100}
\vfill

On fait un \texttt{AND} avec le \alert{complément} du vecteur de Calpurnia (101111).

\vfill
\centerline{Résultat: 100100}

\end{frame}


\begin{frame}{Passons à grande échelle}

Quelques hypothèses:
\begin{itemize}
 \item Un million de documents, mille mots chacun en moyenne.
 \item Disons 6 octets par mot, soit 6 GO (pas très gros en fait!)
 \item Disons 500\,000 termes \alert{distincts}
\end{itemize}

$\Rightarrow$ la matrice a $500 \times 10^9$ bits, soit 62 GO.

\vfill

Ne tient pas en mémoire, ce qui va beaucoup me compliquer les choses....

\vfill

\alert{Mais qu'est-ce qui ne va pas\,?}

\end{frame}


\begin{frame}{La matrice d'incidence est \alert{creuse}}

\begin{block}{Mais oui}
Dans l'ensemble de la matrice, il n'y a que $10^9$ positions avec des 1, soit un sur 500.
\end{block}

Il faut une meilleure structure, appelée \alert{index inversé}.

\figSlideWithSize{inverted-index}{11cm}

On ne peut plus se baser sur la \alert{position}, donc on place dans les cellules \alert{l'identifiant}
du document (\textit{docId}).

\end{frame}

\begin{frame}{Index inversé}

La structure utilisée dans \alert{tous}  les moteurs de recherche.

\begin{itemize}
 \item Un \alert{répertoire} contient tous les \alert{termes}.
 \item Une \alert{liste} (inversée) est associée à chaque \alert{terme}, \alert{triée} par docId.
 \item  Chaque élément de la liste est appelé une \alert{entrée}.
\end{itemize}

\vfill

\alert{Concept}\,: la notion de \alert{terme} (\alert{token} en anglais) est différente de celle de ``mot'', voir plus loin.
\vfill

\alert{Vocabulaire}\,: le répertoire est parfois appelé \emph{dictionnaire}\,; les listes sont des \emph{posting list} en anglais;
les entrées sont des \emph{postings}.

\vfill

\alert{Efficacité}\,: le répertoire devrait toujours être en mémoire\,; les listes, autant que possible en mémoire, sinon
fichiers contigus sur le disque.
\end{frame}

\begin{frame}{Traitement d'une recherche}

Recherchons les documents contenant Brutus et César. 

\vfill

On parcourt \alert{séquentiellement} les deux listes.

\smallskip

On compare à chaque étape les docId\,; s'ils sont égaux\,: \alert{trouvé}!

\smallskip

On avance sur la liste du plus petit docId.

\figSlideWithSize{brutus-et-cesar}{10cm}

\vfill

\alert{Un seul parcours suffit}\,: recherche linéaire, parcours séquentiel. Pas mieux.

\vfill

 C'est une recherche dite \alert{Booléenne}\,: pas de classement, résultat exact.
 
 \vfill
 
 Question\,: et si je ne veux pas Calpurnia\,? Comment je fais\,?
 
\end{frame}


\begin{frame}{L'algorithme (pour les conjonctions)}

\figSlideWithSize{algo-merge}{8cm}

Exercice: écrire l'algorithme de \alert{différence} de deux listes.

\end{frame}


\begin{frame}{Que reste-t-il à faire\,?}

C'est juste la base. Voici quelques problématiques importantes pour la recherche
d'information, nécessitant des extensions / adaptations de l'index inversé de base

\begin{description}
 \item[Quels termes\,?] Quels termes indexe-t-on\,? Beaucoup moins facile que ça n'en a l'air...
 \item[Quelles requêtes\,?] De la plus simple (``sac de mots'') à plus 
                 structurée (index struturé, requêtes Boolénnes, recherche de phrases).
   \item[Performance.] Construction, compression, distribution, optimisation des accès,  mises à jour, etc.
  \item[Classement\,?] Si j'ai des millions de documents dans le résultat, je veux les classer. Comment\,?
\end{description}

On reprend pas à pas.

\end{frame}

\section{Traitement syntaxique et linguistique}

\begin{frame}{Pré-traitement du texte}

Avant d'insérer dans l'index, plusieurs phases.

\vfill

Important: identification de quelques méta-données (la langue), prise en compte du contexte
(quels documents pour quelle application). Puis, de manière  générale\,:

\begin{itemize}
\item \alert{Tokenization}\,: découpage du texte en ``mots''
\item \alert{Normalisation}\,: majuscules\,? acronymes\,? apostrophes\,?

     Exemple\,: \textit{Windows} et \textit{window}, U.S.A vs USA, \emph{l'auditeur}  vs \emph{les auditeurs}, etc.
     
     \item \alert{Stemming} (``racinisation''), \alert{lemmatization}
     
     Prendre la racine des mots pour éviter le biais des variations (auditer, auditeur, audition, etc.)
     
     \item  \alert{Stop words}, quels mots garder\,?
     
     Mots très courants peu informatifs (le, un à, de).
\end{itemize}

C'est de l'art et du réglage... Dans ce qui suit\,: introduction / sensibilisation aux problèmes.

\end{frame}

\begin{frame}{Identification de la Langue}

Comment trouver la langue d'un document\,?

\begin{itemize}

\item Méta information (dans l'entête HTTP p.e.): pas fiable du tout.

\item \alert{Par le jeu de caractères}, pas assez courant!

\figSlide{characters}

\pause

\structure{Respectivement}: Coréen, Japonais, Maldives, Malte, Islandais.

\item Par extension: \alert{séquences de caractères fréquents}, ($n$-grams)

\item Par techniques d'apprentissage (\alert{classifiers})

\end{itemize}

Des librairies font ça très bien (e.g., Tika, \url{http://tika.apache.org})
\end{frame}

\begin{frame}{Tokenisation}

\begin{beamerboxesrounded}{Principe}
Séparation du texte en \alert{tokens} (``mots'')
\end{beamerboxesrounded}

\smallskip


Pas du tout aussi facile qu'on le dirait\,!

\begin{itemize}

\item Dans certaines langues (Chinois, Japonais), les mots  \alert{ne sont pas}
séparés par des espaces.

\item Certaines langues s'écrivent de droite à gauche, de haut en bas.

\item Que faire (et de manière \alert{cohérente}) des acronymes, élisions, nombres, unités, URL, email, etc.

\item \alert{Mots composés}\,: les séparer en \emph{tokens} ou les regrouper en un seul\,?

\begin{enumerate}
   \item Anglais\,: \emph{hostname}, \emph{host-name} et \emph{host
name}, ...
  \item Français\,: Le Mans, aujourd'hui, pomme de terre, ...
  \item Allemand\,: \textit{Levensversicherungsgesellschaftsangestellter} (employé d’une
société d’assurance vie)
\end{enumerate}

 Que faire si l'utilisateur
cherche \emph{hostname} et qu'on a normalisé en \emph{host-name}\,?

\end{itemize}


Majuscules, ponctuation\,? Une solution simple est de normaliser (minuscules, pas de ponctuation).

\end{frame}

\begin{frame}{Exemple pour notre petit jeu de données}

\begin{tabular}{>{\color{structure.fg}}ll}
$d_1$&the$_{\alert{1}}$ jaguar$_{\alert{2}}$ is$_{\alert{3}}$ a$_{\alert{4}}$ new$_{\alert{5}}$ world$_{\alert{6}}$ mammal$_{\alert{7}}$ of$_{\alert{8}}$
the$_{\alert{9}}$ felidae$_{\alert{{10}}}$ family$_{\alert{{11}}}$\\
$d_2$&jaguar$_{\alert{1}}$ has$_{\alert{2}}$ designed$_{\alert{3}}$ four$_{\alert{4}}$ new$_{\alert{5}}$ engines$_{\alert{6}}$\\
$d_3$&for$_{\alert{1}}$ jaguar$_{\alert{2}}$ atari$_{\alert{3}}$ was$_{\alert{4}}$ keen$_{\alert{5}}$ to$_{\alert{6}}$ use$_{\alert{7}}$ a$_{\alert{8}}$
68k$_{\alert{9}}$ family$_{\alert{{10}}}$ device$_{\alert{{11}}}$\\
$d_4$&the$_{\alert{1}}$ jacksonville$_{\alert{2}}$ jaguars$_{\alert{3}}$ are$_{\alert{4}}$ a$_{\alert{5}}$ professional$_{\alert{6}}$
us$_{\alert{7}}$ football$_{\alert{8}}$ team$_{\alert{9}}$\\
$d_5$&mac$_{\alert{1}}$ os$_{\alert{2}}$ x$_{\alert{3}}$ jaguar$_{\alert{4}}$ is$_{\alert{5}}$ available$_{\alert{6}}$ at$_{\alert{7}}$ a$_{\alert{8}}$
price$_{\alert{9}}$ of$_{\alert{{10}}}$ us$_{\alert{{11}}}$ \$199$_{\alert{{12}}}$\\&for$_{\alert{{13}}}$
apple's$_{\alert{{14}}}$ new$_{\alert{{15}}}$ family$_{\alert{{16}}}$ pack$_{\alert{{17}}}$\\
$d_6$&one$_{\alert{1}}$ such$_{\alert{2}}$ ruling$_{\alert{3}}$ family$_{\alert{4}}$ to$_{\alert{5}}$ incorporate$_{\alert{6}}$
the$_{\alert{7}}$ jaguar$_{\alert{8}}$ into$_{\alert{9}}$\\&their$_{\alert{{10}}}$ name$_{\alert{{11}}}$ is$_{\alert{{12}}}$
jaguar$_{\alert{{13}}}$ paw$_{\alert{{14}}}$\\
$d_7$&it$_{\alert{1}}$ is$_{\alert{2}}$ a$_{\alert{3}}$ big$_{\alert{4}}$ cat$_{\alert{5}}$
\end{tabular}

\end{frame}


\begin{frame}{Stemming (racine)}
\begin{beamerboxesrounded}{Principe}
\alert{Confondre} toutes les formes d'un même mot, ou de mots apparentés, en
une seule \alert{racine}.
\end{beamerboxesrounded}

\smallskip

\begin{description}

\item[Stemming Morphologique.] Retire les pluriels, marque de genre, conjugaisons, modes, etc.

\begin{itemize}

\item Très dépendant de la langue\,: \emph{geese} pluriel de \emph{goose}, \emph{mice} de \emph{mouse}

\item Difficile à séparer d'une analyse linguistique (``Les poules du couvent couvent'', ``la petite
    brise la glace''\,:  où est le verbe\,?)
\end{itemize}

\item[Stemming lexical] Fondre les termes proches lexicalement\,: ``politique, politicien, police (?)'' ou 
     ``université, universel, univers (?)'' 


\item[Stemming phonétique.]  Correction fautes de frappes, fautes orthographes

\end{description}

\end{frame}

\begin{frame}{Exemple de stemming}
\begin{tabular}{>{\color{structure.fg}}ll}
$d_1$&the$_1$ jaguar$_2$ \alert{be}$_3$ a$_4$ new$_5$ world$_6$ mammal$_7$ of$_8$
the$_9$ felidae$_{10}$ family$_{11}$\\
$d_2$&jaguar$_1$ \alert{have}$_2$ \alert{design}$_3$ four$_4$ new$_5$
\alert{engine}$_6$\\
$d_3$&for$_1$ jaguar$_2$ atari$_3$ \alert{be}$_4$ keen$_5$ to$_6$ use$_7$ a$_8$
68k$_9$ family$_{10}$ device$_{11}$\\
$d_4$&the$_1$ jacksonville$_2$ \alert{jaguar}$_3$ \alert{be}$_4$ a$_5$ professional$_6$ us$_7$ football$_8$ team$_9$\\
$d_5$&mac$_1$ os$_2$ x$_3$ jaguar$_4$ \alert{be}$_5$ available$_6$ at$_7$ a$_8$
price$_9$ of$_{10}$ us$_{11}$ \$199$_{12}$\\&for$_{13}$
\alert{apple}$_{14}$ new$_{15}$ family$_{16}$ pack$_{17}$\\
$d_6$&one$_1$ such$_2$ \alert{rule}$_3$ family$_4$ to$_5$ incorporate$_6$
the$_7$ jaguar$_8$ into$_9$\\&their$_{10}$ name$_{11}$ \alert{be}$_{12}$
jaguar$_{13}$ paw$_{14}$\\
$d_7$&it$_1$ \alert{be}$_2$ a$_3$ big$_4$ cat$_5$
\end{tabular}
\end{frame}

\begin{frame}{Suppression des \textit{Stop Words}}

\begin{beamerboxesrounded}{Principe}
On retire les mots porteurs d'une information faible afin de limiter le stockage.
\end{beamerboxesrounded}

\begin{description}
\item[articles:] \emph{le}, \emph{le}, \emph{ce}, etc.

\item[verbes ``fonctionnels''] \emph{être}, \emph{avoir}, \emph{faire}, etc.

\item[conjunctions:] \emph{that}, \emph{and}, etc.

\item[etc.]
\end{description}

\begin{remark}
Maintenant moins utilisé car (i) espace de stockage peu coûteux et (ii) pose d'autres problèmes (``pomme de terre'',
``Let it be'', ``Stade de France'')
\end{remark}

\end{frame}



\begin{frame}{Autres problèmes, en vrac}


\alert{Majuscules / minuscules}

Lyonnaise des Eaux, Société Générale, etc.

\vfill
 \alert{Acronymes}
 
   CAT = \textit{cat} ou \textit{Caterpillar Inc.}\,?  M.A.A.F ou MAAF ou Mutuelle ... \,?

\vfill

\alert{Dates, chiffres}

   Monday 24, August, 1572 -- 24/08/1572 -- 24 août 1572
    10000 ou 10,000.00 ou 10,000.00
     
     
     \vfill
     
     \alert{Accents, ponctuation}
     
     résumé ou résume ou resume... 
     
 \begin{remark}
 Dans tous les cas, les même règles de transformation s'appliquent aux documents ET à la requête.
 \end{remark}
 \end{frame}


\begin{frame}{Exemple avec suppression des \textit{stop words}}
\begin{tabular}{>{\color{structure.fg}}ll}
$d_1$&jaguar$_2$ new$_5$ world$_6$ mammal$_7$
felidae$_{10}$ family$_{11}$\\
$d_2$&jaguar$_1$ design$_3$ four$_4$ new$_5$ engine$_6$\\
$d_3$&jaguar$_2$ atari$_3$ keen$_5$
68k$_9$ family$_{10}$ device$_{11}$\\
$d_4$&jacksonville$_2$ jaguar$_3$ professional$_6$
us$_7$ football$_8$ team$_9$\\
$d_5$&mac$_1$ os$_2$ x$_3$ jaguar$_4$ available$_6$
price$_9$ us$_{11}$ \$199$_{12}$
apple$_{14}$\\&new$_{15}$ family$_{16}$ pack$_{17}$\\
$d_6$&one$_1$ such$_2$ rule$_3$ family$_4$ incorporate$_6$
jaguar$_8$ their$_{10}$ name$_{11}$
\\&jaguar$_{13}$ paw$_{14}$\\
$d_7$&big$_4$ cat$_5$
\end{tabular}
\end{frame}



\begin{frame}{Au final, notre index inversé}

\begin{tabular}{>{\color{structure.fg}}ll}
family&$d_1$, $d_3$, $d_5$, $d_6$\\
football&$d_4$\\
jaguar&$d_1$, $d_2$, $d_3$, $d_4$, $d_5$, $d_6$\\
new&$d_1$, $d_2$, $d_5$\\
rule&$d_6$\\
us&$d_4$, $d_5$\\
world&$d_1$\\
\ldots
\end{tabular}

Deux remarques importantes, à l'origine \alert{d'optimisations}.
\begin{itemize}
  \item la longueur des listes inverses est \alert{très} variable\,; 
  \item listes constituées d'entrées très homogènes, permettant des \alert{compressions} efficaces.
\end{itemize}
\end{frame}



\section{Requêtes Booléennes, positionnelles et structurées}

\begin{frame}{Requêtes Booléennes}

Les requêtes sont des \alert{expressions Booléennes}, construites
avec les connecteurs habituels.

\begin{description}
\item[AND] intersection
\item[OR] union
\item[AND NOT] différence
\end{description}


\[\text{(\emph{jaguar} AND \emph{new} AND NOT \emph{family}) OR
\emph{cat}}\]

\vfill

Les documents sont vus comme des \alert{ensembles} de mots.

\vfill

La réponse est \alert{exacte}. Dans le cas d'une conjonction, chaque document contient
\alert{tous} les mots.

\vfill

Toujours très utilisé\,: recherche sur ordinateur personnel, recherche de mails.

\end{frame}

\begin{frame}{Optimisation\,: \alert{d'abord} les petites listes}

La requête \texttt{Brutus AND César AND Calpurnia}

\vfill

\figSlideWithSize{optim-petiteliste}{10cm}

\vfil

\begin{itemize}
  \item On charge la plus petite liste en mémoire\,: c'est le résultat par défaut.
  \item Parcours de  la plus petite parmi celles qui restent\,: \textit{matching} avec la liste en mémoire.
   \item Jusqu'à ce que toutes les listes soient parcourues.
\end{itemize}

\end{frame}

\subsection{Requêtes positionnelles}

\begin{frame}{\textit{Phrase queries}}

Au lieu d'une requête ``sacs de mots'', on ajoute des critères sur les \alert{positions
relatives} des mots de la requête.

\begin{itemize}
  \item par des séquences de mots (``phrases''), ex ``Conservatoire National des Arts et Métiers''.
  
     $\Rightarrow$ ``tout l'art dans cette nation est de conserver son métier'' n'est pas une réponse.
     
  \item par des clauses comme \texttt{NEAR}. 
\end{itemize}

\begin{block}{Extension}
Il faut introduire la \alert{position des mots} dans l'index.
\end{block}

\end{frame}

\begin{frame}{Stockage des positions dans l'index}

Dans chaque entrée, avec chaque document\,: on donne la séquence \alert{triée} des
positions du terme dans le document.

\smallskip

\begin{tabular}{>{\color{structure.fg}}ll}
family&$d_1/11$, $d_3/10$, $d_5/16$, $d_6/4$\\
football&$d_4/8$\\
jaguar&$d_1/2$, $d_2/1$, $d_3/2$, $d_4/3$, $d_5/4$, $d_6/8+13$\\
new&$d_1/5$, $d_2/5$, $d_5/15$\\
rule&$d_6/3$\\
us&$d_4/7$, $d_5/11$\\
world&$d_1/6$\\
\ldots
\end{tabular}

\smallskip

Gros impact sur la \alert{taille} de l'index (l'ordre de grandeur devient celui de la collection elle-même)

\smallskip

Algorithme de recherche légèrement étendu (fusion des positions pour un même document).
\end{frame}

\subsection{Requêtes structurées}


\begin{frame}[fragile]{Requêtes structurées}

Les documents textuels ont toujours une structure plus ou moins stable (regarder ce transparent par exemple).

\vfill

On peut en extraire des \alert{champs} (titre, auteur, date) et les indexer séparément.

\vfill

Le langage de requête permet d'exprimer des contraintes sur ces champs

\medskip

\begin{verbatim}
title CONTAINS 'nfe204'
AND author = 'Rigaux'
AND content CONTAINS 'mongodb'
\end{verbatim}

\vfill

Comment\,? On combine dans le répertoire le terme et le champ (ex: "\texttt{title:nfe204}", "\texttt{author:rigaux}" et 
"\texttt{content:mongodb}").
\vfill
Exemple\,: voir le modèle de données de Lucene et son langage d'interrogation.
\end{frame}

%%%%%%%%%%%%%%%%%%%%%
\section{Construction et mise à jour des index inversés}

%%%%%%%%%%%%%%%%%%%%%%%%%%%%%%%%%%%%%%%%%%%%%%%%%%%%%%%
\begin{frame}
\frametitle{La construction repose sur le \alert{tri}}

Il faut impérativement obtenir des listes triées, de longueur arbitraire. 

\vfill

Pour une collection \alert{statique} de documents\,: 

\begin{itemize}
  \item on extrait les triplets $[d, t, p]$, $d$ étant le document, $t$ le terme et $p$ la position de $t$ dans $d$;

   \item on groupe  les paires $[d, p]$ par terme et on trie chaque groupe\,; 
   
   \item le résultat du tri correspond aux listes inversées.
\end{itemize}

\medskip

On met en {\oe}uvre des techniques de tri classiques, en mémoire, ou sur disque.

\medskip

Dans ce qui suit\,: rappel du tri-fusion en mémoire externe, application aux index inversés.
\end{frame}

%%%%%%%%%%%%%%%%%%%%%%%%%%%%%%%%%%%%%%%%%%%%%%%%%%%%%%%
\begin{frame}
\frametitle{Première phase\,: le tri}


On remplit la mémoire, on trie, on vide
dans des \textbf{fragments}, et
on recommence.

\figSlide{phasetri}

Coût\,: une lecture + une écriture du fichier.

\vfill

On obtient une liste de fragments \alert{triés}

\end{frame}

%%%%%%%%%%%%%%%%%%%%%%%%%%%%%%%%%%%%%%%%%%%%%%%%%%%%%%%
\begin{frame}
\frametitle{Deuxième phase\,: la fusion}


On groupe les fragments par $M$ (taille de la zone
de tri), et on fusionne.\\

\figSlideWithSize{trifusion}{9cm}

Coût\,: autant de lectures/écritures du fichier que de niveaux
de fusion.

\end{frame}

%%%%%%%%%%%%%%%%%%%%%%%%%%%%%%%%%%%%%%%%%%%%%%%%%%%%%%%
\begin{frame}
\frametitle{Illustration avec $M = 3$}

\figSlideWithSize{tri-fusion}{10cm}

\end{frame}
%%%%%%%%%%%%%%%%%%%%%%%%%%%%%%%%%%%%%%%%%%%%%%%%%%%%%%%
\begin{frame}

\frametitle{Essentiel\,: la taille de la zone de tri}

Un fichier de 75~000 pages de 4 Ko, soit 307 Mo.

\vfill

 \textbf{$M > 307 Mo$}\,: une lecture, soit 307 

\vfill

 \textbf{$M = 2 Mo$}, soit 500 pages.
    \begin{enumerate}
      \item le tri donne $\lceil \frac{307}{2} \rceil = 154$
        fragments.

      \item On fait la fusion avec $154$ pages
     \end{enumerate}
       Coût total de $614 + 307 = 921$ Mo. 

\vfill

NB\,: il faut allouer beaucoup de mémoire pour
passer de 1 à 0 niveau de tri.
\end{frame}

%%%%%%%%%%%%%%%%%%%%%%%%%%%%%%%%%%%%%%%%%%%%%%%%%%%%%%%
\begin{frame}

\frametitle{Avec très peu de mémoire}
\textbf{$M = 1 Mo$}, soit 250 pages.
    \begin{enumerate}
       \item on obtient 307 fragments.

       \item On fusionne les 249 premiers fragments,
               puis les 58 restant. On obtient $F_1$ et $F_2$.
        \item On fusionne $F_1$ et $F_2$.
   \end{enumerate}

\medskip

 Coût total\,: $1\,228 + 307 = 1\,535$ Mo.

\medskip

Résultat\,: grosse dégradation entre 2 Mo et 1 Mo (calcul
approximatif).

\end{frame}


%%%%%%%%%%%%%%%%%%%%%%%%%%%%%%%%%%%%%%%%%%%%%%%%%%%%%%%
\begin{frame}
\frametitle{Construction d'un index inversé}

On crée un premier index en mémoire, puis on l'écrit sur  disque.

\figSlideWithSize{merge-based}{10cm}

On \alert{fusionne} périodiquement avec les mises à jour suivantes.

\begin{remark}
Pas adapté à des mises à jour intensives.
\end{remark}
\end{frame}


\begin{frame}{Attention à la taille de l'index}
 
Avec les positions (et d'autres informations éventuelles, comme le poids, cf. plus loin), la taille de l'index
devient un problème.

\smallskip

Supposons qu'on indexe une collection d'un million d'emails, de taille moyenne 1,000 octets. Chaque
email contient 100 mots, et il y a 200\,000 termes distincts.


\begin{enumerate}
  \item taille de la collection, nombre d'occurrences de mots\,? % 1GBn 100 millions
  \item combien de listes dans l'index\,? % 200,000 listes;
  \item supposons que  20\% des termes d'un document apparaissent deux fois\,; dans combien de listes
      un document apparaît-il\,?  % 80 lists, 
  \item combien d'entrées dans une liste\,? %  500 entries;
  \item  un docId et une position sont codée sur 4 octets, quelle est la taille de l'index\,? % 800 MO
  
%the \emph{average} size of a list is 1,600 bytes; the whole index contains
% $400 \times 200{,}000 = 80{,}000{,}000$ entries; the index size is 320 MB (for inverted lists) plus 
%   2,4 MB ($200{,}000 \times 12$) for the directory. 
\end{enumerate}
\end{frame}


\begin{frame}{Compression des listes inversées}

Sans compression, la taille d'un index peut devenir plus importante que la collection elle-même.

\vfill
Une bonne technique de compression est essentielle. 

\vfill

Note\,: le temps de décompression doit représenter un gain par rapport à la lecture de la liste non-compressée.

\vfill
On s'appuie sur le fait que les docId et les positions sont stockés en ordre ascendant.

\[[87; 273; 365; 576; 810].\]

\vfill

Première étape, on utilise le \emph{delta-coding}:

\[[87; 186; 92; 211; 234].\]

Exercice: quel est le nombre minimal d'octets pour représenter la première liste\,? la seconde\,?
\end{frame}


\begin{frame}{Codage par octets variables}

Idée: on code les entiers sur 7 bits ($2^7= 128$); le premier bit à 1 indique la terminaison de l'encodage du nombre.

\vfill

Soit  $v=9$, on l'encode sur un octet $10000101$ (premier bit à 1). 

\vfill
Soit $v=137$. 
\begin{enumerate}
  \item le premier octet encode $v'= v \bmod 128 = 9$, donc $b=10000101$ comme précédemment;
  \item ensuite on encode $ v / 128 = 1 $, avec un octet $b' = 00000001$ (premier bit à 0).
\end{enumerate}

$137$ est donc encodé sur deux octets:
\[00000001\,\,10000101.\]
\vfill
Taux de compression\,: typiquement 1/4 à 1/2 de la représentation fixe.
\end{frame}


\end{document}

\section{Requêtes avec classement}

\begin{frame}{Ranked search}

Boolean search does not give an accurate result because it does not take account
of the \alert{relevance} of a document to a query.

\vfill

If the search retrieves dozen or hundreds of documents, the most relevant must
be shown in top position!

\vfill

The quality of a result with respect to relevance is measured by two factors:

$$precision =\frac{|relevant| \cap |retrieved|}{|retrieved|}$$

$$recall =\frac{|relevant| \cap |retrieved|}{|relevant|}$$

\end{frame}


\begin{frame}{Weighting terms occurrences}

Relevance can be computed by giving a \alert{weight}
to term occurrences.
\begin{itemize}

\item Terms occurring \alert{frequently} in a \alert{given document}:
more \textcolor{magenta}{relevant}.


The \emph{term frequency} is the number of occurrences of a term $t$ in a
document $d$, divided by the total number of terms in $d$ (normalization)

\[\mathop{\mathrm{tf}}(t,d)=\frac{n_{t,d}}{\sum_{t'} n_{t',d}}\]
\noindent
where $n_{t',d}$ is the number of occurrences of $t'$ in $d$.  

\item Terms occurring \alert{rarely} in the \alert{document collection}
as a whole: more \textcolor{orange}{informative}

The \emph{inverse document frequency}  (idf) is obtained from the division of 
the total number of documents by the number of documents where $t$ occurs, as follows:
\[\mathop{\mathrm{idf}}(t)=\log\frac{|D|}{\left|\left\{d'\in
D\,|\,n_{t,d'}>0\right\}\right|}.\]

\end{itemize}

\end{frame}

\begin{frame}{TF-IDF Weighting}

The inverted is extended by adding 
\alert{Term Frequency---Inverse Document Frequency} weighting
 \[
\mathop{\mathsf{tf{}idf}}(t,d)={\color{magenta}\frac{n_{t,d}}{\sum_{t'}
n_{t',d}}}\cdot
\color{orange}\log\frac{|D|}{\left|\left\{d'\in D\,|\,n_{t,d'}>0\right\}\right|}\]
\begin{center}\begin{tabular}{cl}
$n_{t,d}$&number of occurrences of $t$ in $d$\\
$D$&set of all documents
\end{tabular}
\end{center}

Documents (along with weight) are stored in \alert{decreasing weight
order} in
the index

\end{frame}


\begin{frame}{TF-IDF Weighting Example}
\begin{tabular}{>{\color{structure.fg}}ll}
family&$d_1/11/.13$, $d_3/10/.13$, $d_6/4/.08$, $d_5/16/.07$\\
football&$d_4/8/.47$\\
jaguar&$d_1/2/.04$, $d_2/1/.04$, $d_3/2/.04$, $d_4/3/.04$,
  $d_6/8+13/.04$,\\& $d_5/4/.02$\\
new&$d_2/5/.24$, $d_1/5/.20$, $d_5/15/.10$\\
rule&$d_6/3/.28$\\
us&$d_4/7/.30$, $d_5/11/.15$\\
world&$d_1/6/.47$\\
\ldots
\end{tabular}

Exercise: take an entry, and check that the tf/idf value is indeed correct
(take documents after stop-word removal).

\end{frame}


\begin{frame}{Answering Top-$k$ Queries}

 \[t_1\text{ AND }\dots\text{ AND }t_n\]

\begin{beamerboxesrounded}{Problem}
Find the \alert{top-$k$ results} (for some given $k$) 
to the query, without retrieving all documents
matching it.
\end{beamerboxesrounded}

\vspace{.5cm}

Notations:
\begin{description}
\item[$s(t,d)$] weight of $t$ in $d$ (e.g., $\mathsf{tf{}idf}$)
\item[$g(s_1,\dots,s_n)$] monotonous function that computes the global
score (e.g., addition)
\end{description}
\end{frame}



\begin{frame}[fragile]
\frametitle{Exercise}
Consider the following documents:
\begin{enumerate}
  \item $d_1$ = I like to watch the sun set with my friend.
  \item $d_2$ = The Best Places To Watch The Sunset.
  \item $d_3$ = My friend watches the sun come up.
\end{enumerate}

Construct an inverted index with tf/idf weights for terms 'Best' and 'sun'. What
would be the ranked result of the query 'Best and sun'?

\end{frame}

\begin{frame}{Basic algorithm}

First version of the top-$k$ algorithm: the inverted file contains entries
sorted on the document id. The query is \[t_1\text{ AND }\dots\text{ AND }t_n\].

\begin{enumerate}
\item Take the first entry of each list; one obtains a tuple $T = [e_1, \ldots
e_n]$;
\item Let $d_1$ be the minimal doc id in the entries of $T$: compute the global
score of $d_1$;
 \item For each entry $e_i$ featuring $d_1$: advance on the inverted list $L_i$.
\end{enumerate}

When \emph{all} lists have been scanned: sort the documents on the global
scores.
\vfill

Not very efficient; cannot give the ranked result before a full scan on the
lists.
\end{frame}

\begin{frame}{The Threshold Algorithm}
\begin{enumerate}
\item Let $R$ be the empty list, and $m=+\infty$.
\item For each $1\leq i\leq n$:
\label{step:main-fagin}
\begin{enumerate}
\item Retrieve the document $d^{(i)}$ containing term
$t_i$ that has the \alert{next largest} $s(t_i,d^{(i)})$.
\item Compute its global score
$g_{d^{(i)}}=g(s(t_1,d^{(i)}),\dots,s(t_n,d^{(i)}))$
by retrieving all $s(t_j,d^{(i)})$ with $j\neq i$.
\item If $R$ contains less
than $k$ documents, or if $g_{d^{(i)}}$ is greater than the \alert{minimum of
the score of
documents} in $R$, add $d^{(i)}$ to $R$.
\end{enumerate}
\item Let $m=g(s(t_1,d^{(1)}),s(t_2,d^{(2)}),\dots,s(t_n,d^{(n)}))$.
\item If $R$
contains more than $k$ documents, and the minimum of the score of the
documents in $R$ is \alert{greater than or equal} to $m$, return $R$.
\item Redo step \ref{step:main-fagin}.
\end{enumerate}
\end{frame}

\begin{frame}{The TA, by example}

$q$ = ``\emph{new} OR \emph{family}'', and $k=3$. 
We use inverted
lists sorted on the weight. 

\begin{tabular}{>{\bfseries}p{4em}l}
family&$d_1/11/.13$, $d_3/10/.13$, $d_6/4/.08$, $d_5/16/.07$\\
new&$d_2/5/.24$, $d_1/5/.20$, $d_5/15/.10$\\
\ldots
\end{tabular}

Initially, $R=\varnothing$ and $\tau=+\infty$. 

\begin{enumerate}
  \item $d^{(1)}$ is the first entry in $L_{\rm{family}}$, one finds 
  $s(\rm{new}, d_1) = .20$; the global score
for $d_1$ is $.13 + .20 = .33$.

\item Next, $i=2$, and one finds that the global score for $d_2$ is .24.

\item The
algorithm quits the loop on $i$ with $R = \langle[d_1,.33], [d_2, .24] \rangle$ 
and $\tau = .13 + .24 = .37$.

\item We proceed with the loop
again, taking $d_3$ with score $.13$ and $d_5$ with score $.17$. 
$[d_5, .17]$ is added to $R$ (at the end) and  $\tau$ is now $.10 + .13 = .23$.

  A last loop concludes that the next candidate
is $d_6$, with a global score of $.08$. Then we are done. 

\end{enumerate}

\end{frame}


\end{document}



\subsection{Clustering}

\begin{frame}{Clustering Example}
\includegraphics[width=\linewidth,trim={0 500 0 0},clip]{../figures/clusty}
\end{frame}

\begin{frame}{Cosine Similarity of Documents}

\begin{itemize}
\item \alert{Document Vector Space} model:
\begin{description}
\item[terms] dimensions
\item[documents] vectors
\item[coordinates] weights
\end{description}

(The projection of document $d$ along coordinate $t$ is the weight of $t$
in~$d$, say $\mathsf{tf{}idf(t,d)}$)

\item Similarity between documents $d$ and $d'$: \alert{cosine} of these two
vectors

\[
\alert{\cos(d,d')}=\frac{d\cdot d'}{\Vert d\Vert\times\Vert d'\Vert}
\]
\begin{center}\begin{tabular}{cl}
$d\cdot d'$&scalar product of $d$ and $d'$\\
$\Vert d\Vert$&norm of vector $d$
\end{tabular}\end{center}

\item $\cos(d,d)=1$

\item $\cos(d,d')=0$ if $d$ and $d'$ are \alert{orthogonal} (do not
share any common term)
\end{itemize}
\end{frame}

\begin{frame}{Agglomerative Clustering of Documents}

\begin{enumerate}
\item Initially, each document forms \alert{its own cluster}.
\item The similarity between two clusters is defined as the \alert{maximal
similarity} between elements of each cluster.
\item Find the two clusters whose mutual similarity is \alert{highest}. If it is
\alert{lower than a given threshold}, end the clustering. Otherwise, regroup
these clusters. Repeat.
\end{enumerate}

\begin{remark}
Many other more refined algorithms for clustering exist.
\end{remark}

\end{frame}

\subsection{Other Media}

\begin{frame}{Indexing HTML}
\begin{itemize}
\item HTML: text + \alert{meta-information} + \alert{structure}

\item Possibly: \alert{separate index} for meta-information (title, keywords)

\item \alert{Increase weight} of structurally emphasized content in index

\item \alert{Tree structure} can also be queried with XPath or XQuery, but not
very useful on the Web as a whole, because of tag soup and lack of
consistency.
\end{itemize}
\end{frame}

\begin{frame}{Indexing Multimedia Content}

\begin{itemize}

\item Basic approach: index \alert{text from context} of the media
\begin{itemize}
\item surrounding text
\item text in or around the links pointing to the content
\item filenames
\item associated subtitles (hearing-impaired track on TV)
\end{itemize}

\item Elaborate approach: index and search the media itself, with the
help of \alert{speech recognition} and \alert{sound, image, and video
analysis}
\begin{itemize}
\item TrackID, Shazam: identify a song played on the radio
\item Musipedia: look for music by whistling a tune,
\url{http://www.musipedia.org/}, \url{http://www.midomi.com/}
\item Image search from a similar image,
\url{http://images.google.com/}, Google Goggles
\item Voxalead, \url{http://voxaleadnews.labs.exalead.com/}
\end{itemize}
\end{itemize}

\end{frame}

\end{document}
